The results in this section follow the corresponding topics in \cite{burton}.
The checked declarations belong to \texttt{ElementaryNumberTheory.Chapter04}.

\begin{theorem}
  \label{nt:thm:wilson}
  If $p$ is prime, then $(p-1)!
    \equiv-1\pmod p$.

  \begin{lean}
    theorem wilson (p : ℕ) (hp : Chapter02.isPrime p) :
        ((p - 1).factorial : ℤ) ≡ -1 [ZMOD (p : ℤ)]
  \end{lean}
\end{theorem}

\begin{proof}
  Pair every nonzero residue with its inverse.
  A fixed point satisfies $x^2=1$, hence is one or minus one.
  After removing minus one, all inverse pairs multiply to one, and the possible fixed point one
  does too.
  The remaining factor is minus one.
  The code identifies the product of the nonzero residues with the factorial, using the
  inverse-pair lemma proved in the module.

  \begin{lean}
    classical

    have : Fact p.Prime := ⟨Nat.prime_def.mpr hp⟩

    have hpos : 0 < p := by
      have := hp.1
      omega

    have hprod : (∏ u : (ZMod p)ˣ, (u : ZMod p)) =
        ∏ k ∈ Finset.Ico 1 p, (k : ZMod p) := by

      apply Finset.prod_bij (fun (u : (ZMod p)ˣ) _ => (u : ZMod p).val)

      · intro u hu
        apply Finset.mem_Ico.mpr
        refine ⟨?_, ZMod.val_lt _⟩

        have hne : (u : ZMod p).val ≠ 0 := by
          intro hz
          apply Units.ne_zero u
          apply ZMod.val_injective p
          simpa only [ZMod.val_zero] using hz
        omega

      · intro u hu v hv h
        exact Units.ext (ZMod.val_injective p h)

      · intro k hk

        have hk' := Finset.mem_Ico.mp hk

        have hne : (k : ZMod p) ≠ 0 := by
          intro hz

          have hh := congrArg ZMod.val hz
          rw [ZMod.val_natCast_of_lt hk'.2, ZMod.val_zero] at hh
          omega

        exact ⟨Units.mk0 (k : ZMod p) hne, Finset.mem_univ _,
          ZMod.val_natCast_of_lt hk'.2⟩

      · intro u hu
        exact (ZMod.natCast_zmod_val _).symm

    have hf : (∏ k ∈ Finset.Ico 1 p, (k : ZMod p)) = ((p - 1).factorial : ZMod p) := by
      rw [← Nat.cast_prod, ← Finset.prod_Ico_id_eq_factorial]
      rw [Nat.sub_add_cancel hpos]

    apply (ZMod.intCast_eq_intCast_iff _ _ p).1
    simpa only [Int.cast_natCast, Int.cast_neg, Int.cast_one] using
      (hf.symm.trans (hprod.symm.trans (prod_units_eq_neg_one p)))
  \end{lean}
\end{proof}

\begin{theorem}
  \label{nt:thm:prime-of-wilson}
  If $n>1$ and $(n-1)!
    \equiv-1\pmod n$, then $n$ is prime.

  \begin{lean}
    theorem prime_of_wilson (n : ℕ) (hn : 1 < n)
        (h : ((n - 1).factorial : ℤ) ≡ -1 [ZMOD (n : ℤ)]) : Chapter02.isPrime n
  \end{lean}
\end{theorem}

\begin{proof}
  Let $d$ divide $n$.
  If $d\ne n$, then $1\le d\le n-1$, so $d$ divides $(n-1)!
  $.
  It also divides $(n-1)!
    +1$ by the congruence, and therefore divides one.
  Thus every divisor is one or $n$.

  \begin{lean}
    refine ⟨hn, ?_⟩
    intro d hd
    by_cases he : d = n

    · exact Or.inr he

    · have hdpos : 0 < d := by
        obtain ⟨k, hk⟩ := hd
        by_contra hh

        have hz : d = 0 := by omega
        rw [hz, zero_mul] at hk
        omega

      have hdle : d ≤ n := by
        obtain ⟨k, hk⟩ := hd

        have hkpos : 0 < k := by nlinarith
        nlinarith

      have hdf : (d : ℤ) ∣ (n - 1).factorial :=
        Int.natCast_dvd_natCast.mpr (dvd_factorial d (n - 1) hdpos (by omega))

      obtain ⟨u, hu⟩ := hdf
      obtain ⟨v, hv⟩ := Int.natCast_dvd_natCast.mpr hd
      obtain ⟨k, hk⟩ := (Chapter03.modEq_iff_dvd_sub _ _ _).1 h

      have hone : (1 : ℤ) = d * (v * k - u) := by
        rw [hu, hv] at hk
        nlinarith

      have hdz : (0 : ℤ) < d := by exact_mod_cast hdpos

      have hkp : 0 < v * k - u := by nlinarith
      left

      have : (d : ℤ) = 1 := by nlinarith
      exact_mod_cast this
  \end{lean}
\end{proof}

\begin{theorem}
  \label{nt:thm:prime-iff-wilson}
  For $n>1$, primality is equivalent to $(n-1)!
    \equiv-1\pmod n$.

  \begin{lean}
    theorem prime_iff_wilson (n : ℕ) (hn : 1 < n) :
        Chapter02.isPrime n ↔ ((n - 1).factorial : ℤ) ≡ -1 [ZMOD (n : ℤ)]
  \end{lean}
\end{theorem}

\begin{proof}
  Combine Wilson\textquotesingle s theorem with its converse.

  \begin{lean}
    ⟨wilson n, prime_of_wilson n hn⟩
  \end{lean}
\end{proof}

\begin{definition}
  A quadratic congruence modulo $n$ has coefficients $a,b,c\in\mathbb Z$ and the form
  $ax^2+bx+c\equiv0\pmod n$, with $a\not\equiv0\pmod n$.

  \begin{lean}
    structure QuadraticCongruence (n : ℤ) where
      a : ℤ
      b : ℤ
      c : ℤ
      leading_ne_zero : ¬a ≡ 0 [ZMOD n]
  \end{lean}
\end{definition}

\begin{theorem}
  \label{nt:thm:half-factorial}
  If $p=2m+1$, then in the residue ring modulo $p$, \[ (p-1)!=(-1)^m(m!)^2. \]

  \begin{lean}
    theorem half_factorial (p m : ℕ) (hp : p = 2 * m + 1) :
        ((p - 1).factorial : ZMod p) = (-1 : ZMod p) ^ m * (m.factorial : ZMod p) ^ 2
  \end{lean}
\end{theorem}

\begin{proof}
  Split the factorial into the factors $1,\ldots,m$ and $m+1,\ldots,2m$.
  Reflection $k\mapsto p-k$ carries the second interval to the first, and each reflected residue is
  $-k$.

  \begin{lean}
    have hprefix : (∏ k ∈ Finset.Ico 1 (m + 1), (k : ZMod p)) = (m.factorial : ZMod p) := by
      rw [← Nat.cast_prod, Finset.prod_Ico_id_eq_factorial]

    have hreflect : (∏ k ∈ Finset.Ico (m + 1) p, (k : ZMod p)) =
        ∏ k ∈ Finset.Ico 1 (m + 1), ((p - k : ℕ) : ZMod p) := by
      rw [Finset.prod_Ico_reflect _ 1 (by omega : m + 1 ≤ p + 1)]
      rw [show p + 1 - (m + 1) = m + 1 by omega, Nat.add_sub_cancel]

    have hnegative : (∏ k ∈ Finset.Ico 1 (m + 1), ((p - k : ℕ) : ZMod p)) =
        (-1 : ZMod p) ^ m * (m.factorial : ZMod p) := by
      have hterm : ∀ k ∈ Finset.Ico 1 (m + 1), ((p - k : ℕ) : ZMod p) = -1 * (k : ZMod p) := by
        intro k hk

        have hkp : k ≤ p := by
          have := Finset.mem_Ico.mp hk
          omega

        rw [Nat.cast_sub hkp, ZMod.natCast_self, zero_sub, neg_one_mul]

      rw [Finset.prod_congr rfl hterm, Finset.prod_mul_distrib, Finset.prod_const,
        Nat.card_Ico, Nat.add_sub_cancel, hprefix]
  \end{lean}

  The first product is $m!
  $ and the second is $(-1)^m m!$.
  Multiplying the two pieces gives the identity.

  \begin{lean}
    have hfull : ((p - 1).factorial : ZMod p) = ∏ k ∈ Finset.Ico 1 p, (k : ZMod p) := by
      rw [← Nat.cast_prod, ← Finset.prod_Ico_id_eq_factorial]
      rw [Nat.sub_add_cancel (by omega : 1 ≤ p)]

    rw [hfull, ← Finset.prod_Ico_consecutive (fun k => (k : ZMod p))
      (by omega : 1 ≤ m + 1) (by omega : m + 1 ≤ p), hprefix, hreflect, hnegative]
    ring
  \end{lean}
\end{proof}

\begin{theorem}
  \label{nt:thm:exists-square-eq-neg-one-iff}
  For an odd prime $p$, the congruence $x^2\equiv-1\pmod p$ has an integer solution exactly when
  $p\equiv1\pmod4$.

  \begin{lean}
    theorem exists_square_eq_neg_one_iff (p : ℕ) (hp : Chapter02.isPrime p) (hodd : Odd p) :
        (∃ x : ℤ, x ^ 2 ≡ -1 [ZMOD (p : ℤ)]) ↔ p % 4 = 1
  \end{lean}
\end{theorem}

\begin{proof}
  Write $p=2m+1$.
  If $x^2\equiv-1$, then $p\nmid x$.
  Raising the square congruence to the $m$th power and using Fermat gives $(-1)^m\equiv1$.
  If $p\equiv3\pmod4$, then $m$ is odd, forcing $p\mid2$, contrary to $p>2$.

  \begin{lean}
    have hp2 : 2 < p := by
      have hmod := Nat.odd_iff.mp hodd

      have := hp.1
      omega

    let m := (p - 1) / 2

    have hpm : p = 2 * m + 1 := by
      have hmod := Nat.odd_iff.mp hodd
      dsimp [m]
      omega

    constructor

    · rintro ⟨x, hx⟩

      have hnot : ¬(p : ℤ) ∣ x := by
        intro hd

        have hz := Chapter03.modEq_pow ((Chapter03.modEq_zero_iff _ _).2 hd) 2
        rw [zero_pow (by decide : 2 ≠ 0)] at hz

        have hh := (Chapter03.modEq_iff_dvd_sub _ _ _).1
          (Chapter03.modEq_trans (Chapter03.modEq_symm hz) hx)

        obtain ⟨k, hk⟩ := hh

        have hpz : (2 : ℤ) < p := by exact_mod_cast hp2

        have hkpos : 0 < k := by nlinarith
        nlinarith

      have hpow := Chapter03.modEq_pow hx m
      rw [← pow_mul, show 2 * m = p - 1 by omega] at hpow

      have hh := Chapter03.modEq_trans (Chapter03.modEq_symm hpow) (fermat p hp x hnot)
      by_contra hne

      have hm : Odd m := Nat.odd_iff.mpr (by omega)
      rw [hm.neg_one_pow] at hh

      obtain ⟨k, hk⟩ := (Chapter03.modEq_iff_dvd_sub _ _ _).1 hh

      have hpz : (2 : ℤ) < p := by exact_mod_cast hp2

      have hkneg : k < 0 := by nlinarith
      nlinarith
  \end{lean}

  If $p\equiv1\pmod4$, then $m$ is even.
  The half-factorial identity and Wilson\textquotesingle s theorem give $(m!
    )^2\equiv-1$, which supplies the required solution.

  \begin{lean}
    · intro hpmod

      have hm : Even m := Nat.even_iff.mpr (by omega)

      have hhalf := half_factorial p m hpm
      rw [hm.neg_one_pow, one_mul] at hhalf

      have hw := (ZMod.intCast_eq_intCast_iff _ _ p).2 (wilson p hp)
      simp only [Int.cast_natCast, Int.cast_neg, Int.cast_one] at hw

      refine ⟨m.factorial, ?_⟩
      apply (ZMod.intCast_eq_intCast_iff _ _ p).1
      simpa only [Int.cast_pow, Int.cast_natCast, Int.cast_neg, Int.cast_one] using
        hhalf.symm.trans hw
  \end{lean}
\end{proof}

\begin{theorem}
  \label{nt:thm:infinite-factorial-composites}
  Infinitely many composite integers have the form $n!
    +1$.

  \begin{lean}
    theorem infinite_factorial_composites :
        {k : ℕ | (∃ n : ℕ, k = n.factorial + 1) ∧ Chapter02.isComposite k}.Infinite
  \end{lean}
\end{theorem}

\begin{proof}
  Choose primes $p$ arbitrarily large and put $n=p-1$.
  Wilson gives $p\mid n!
    +1$.
  For $n\ge3$, the elementary bound $n!
    \ge2n$ makes this divisor proper.
  These composite values exceed every bound.

  \begin{lean}
    apply Set.infinite_of_forall_exists_gt
    intro B

    obtain ⟨n, hB, hn⟩ := exists_composite_factorial_add_one_gt B

    exact ⟨n.factorial + 1, ⟨⟨n, rfl⟩, hn⟩, hB⟩
  \end{lean}
\end{proof}
